\subsection{Within-Country Comparison}\label{sec:methods_data}

Per analizzare l’home advantage è stato effettuato un confronto within-country: la performance di ogni Paese che ha ospitato almeno una volta le Olimpiadi è stata confrontata con la media delle performance dello stesso Paese quando non ha ospitato.

Dal dataset originale sono quindi state estratte le righe relative a tutti i Paesi che hanno ospitato almeno una volta le Olimpiadi.

Fatto ciò, è stata calcolata la media delle percentuali di medaglie vinte da ogni Paese nelle edizioni in cui i paesi erano ospitanti e la media delle percentuali di medaglie vinte quando invece non lo erano. Sono stati poi calcolati i delta, ovvero le differenze fra queste due misure.

\subsection*{Sign-Test}
Il Test dei Segni è un test statistico non parametrico utilizzato per determinare se esiste una differenza sistematica tra coppie di osservazioni correlate. Viene spesso impiegato come alternativa al test t di Student o al test di Wilcoxon quando non è possibile garantire le assunzioni sulla distribuzione dei dati.

In particolare, il test t di Student richiede che la distribuzione dei dati sia normale, mentre il test di Wilcoxon richiede che le differenze che si vogliono analizzare siano distribuite in maniera simmetrica.

il Test dei Segni si basa esclusivamente sulla direzione della differenza tra le coppie. Per ogni coppia di dati (x, y), il test assegna:

\begin{itemize}
    \item Un segno più se x $>$ y;
    \item Un segno meno se x $<$ y;
    \item I casi di parità (differenza pari a zero) vengono solitamente esclusi dall'analisi.
\end{itemize}

Sotto l'ipotesi nulla H0, ci si aspetta che il numero di segni positivi e negativi sia approssimativamente uguale, seguendo una distribuzione binomiale con probabilità p-value pari a 0.5.

Sebbene sia molto flessibile, il Test dei Segni ha una minore potenza statistica rispetto al test di Wilcoxon, poiché ignora l'entità numerica delle differenze. Tuttavia, in presenza di distribuzioni fortemente non simmetriche, la sua affidabilità lo rende preferibile per evitare errori di tipo I (falsi positivi).

Nel nostro caso, non conoscendo la distribuzione dei dati e avendo verificato una forte asimmetria, non è stato possibile usare il test t di Student o il test di Wilcoxon. Le differenze analizzate dal test dei segni sono i delta calcolati nel Within-Country Comparison.

\subsection*{Spearman Correlation}
Ci sono fattori sistematici che influenzano la performance dei vari Paesi alle Olimpiadi?

Nella letteratura esistente, due dei fattori che più vengono tenuti in considerazione sono la ricchezza e la popolazione di un Paese:
\begin{itemize}
    \item la ricchezza, in quanto permette al Paese di investire nelle infrastrutture sportive e nel sostegno agli atleti
    \item la popolazione, in quanto al crescere del numero di abitanti, probabilisticamente, cresce la dimensione del bacino di potenziali talenti
\end{itemize}

Per verificare correlazione tra ricchezza/popolazione e performance Olimpica, abbiamo deciso di utilizzare la correlazione di Spearman.

Il coefficiente di correlazione di Spearman è una misura statistica non parametrica che valuta la forza e la direzione della relazione tra due variabili. A differenza della correlazione di Pearson, che misura la relazione lineare, Spearman valuta la relazione monotonica.

Il calcolo non si basa sui valori grezzi dei dati, ma sui loro ranghi (ovvero sulla loro posizione in classifica). Il procedimento prevede di:

\begin{enumerate}
    \item Ordinare i dati di ciascuna variabile dal più piccolo al più grande.
    \item Assegnare un rango (1, 2, 3...) a ogni valore.
    \item Calcolare il coefficiente di correlazione basandosi sulla differenza tra i ranghi delle due variabili.
\end{enumerate}

Abbiamo deciso di utilizzare Spearman invece del più classico Pearson per due ragioni principali:
\begin{itemize}
    \item Spearman non richiede che i dati seguano una distribuzione normale
    \item Spearman è in grado di catturare relazioni che Pearson ignorerebbe, come ad esempio una crescita esponenziale o logaritmica. Spearman non richiede che i dati crescano secondo una linea retta; è sufficiente che, al crescere di una variabile, l'altra tenda a crescere (o decrescere) in modo costante, anche se con un ritmo non uniforme.
\end{itemize}

Il coefficiente di Spearman varia da -1 (perfetta relazione monotonica negativa) a +1 (perfetta relazione monotonica positiva). Uno 0 indica assenza di relazione.
